\documentclass[10pt,a4paper]{amsart}
\usepackage[margin=19mm]{geometry}
\usepackage{amsmath,amssymb,mathtools}
\usepackage[scr=rsfs,cal=euler]{mathalfa}
\usepackage{xfrac}
\usepackage[unicode,hidelinks,bookmarks=false]{hyperref}
\newcommand{\ie}{i.e.\ }
\newcommand{\cf}{cf.\ }
\newcommand{\resp}{respectively\ }
\newcommand{\Subsection}{\subsection}
\providecommand{\nobreakdash}{\nobreak\hbox{-}}
\setlength{\emergencystretch}{2em}
\multlinegap0pt
\usepackage{mathrsfs}
\usepackage{booktabs}
\usepackage{array,enumitem}
\usepackage{tikz}
\usetikzlibrary{arrows.meta, math,positioning, calc}
\newtheorem{theorem}{Theorem}[section]
\newtheorem{proposition}[theorem]{Proposition}
\newtheorem{lemma}[theorem]{Lemma}
\newtheorem{corollary}[theorem]{Corollary}
\newtheorem{claim}[theorem]{Claim}
\newtheorem*{theorem*}{Theorem}
\newtheorem{mainthm}{Theorem}
\renewcommand{\themainthm}{\Alph{mainthm}}
\newtheorem{definition}[theorem]{Definition}
\newtheorem{principle}[theorem]{Principle}
\newtheorem{example}[theorem]{Example}
\newtheorem{remark}[theorem]{Remark}
\newtheorem{question}[theorem]{Question}
\newtheorem{conjecture}[theorem]{Conjecture}
\newtheorem{observation}[theorem]{Observation}
\numberwithin{equation}{section}
\def\C{{\mathbb{C}}}
\def\D{\mathcal{D}}
\def\M{\mathscr{M}}
\def\Com#1{\mathbb{C}^{#1}}
\def\R#1{\mathbb{R}^{#1}}
\def\Z#1{\mathbb{Z}^{#1}}
\def\Tcirc{\mathbb{T}}
\def\Haus#1{\mathcal{H}^{#1}}
\def \const{\mathrm{const}}
\def\Del{\frak{I}}
\def\Dd#1{\mathbb{D}\left(#1\right)}
\def\H{\mathcal{H}}
\def\DD{\mathbb{D}}
\def\llangle{\langle\!\langle}
\def\rrangle{\rangle\!\rangle}
\def\ZN{\mathbb{Z}_N^\circledast}
\def\QLQ{\quad \Leftrightarrow\quad}
\def\con#1{\overline{#1}}
\def\z{\zeta}
\def\x{\overline{x}}
\def\ov#1{\overline{#1}}
\def\fr#1{\frac{1}{#1}}
\def\scal#1#2{\langle #1; #2 \rangle}
\def\scalh#1#2{\scal{#1}{#2}_{\frak{H}(t)}}
\def\vv{\mathbf{v}}
\def\conj#1{\overline{#1}}
\def\V{\frak{V}}
\def\scalm#1#2{\scal{#1}{#2}_{\V}}
\def\scalM#1#2{\{#1;#2\}}
\def\diff#1#2{\frac{d#1}{d#2}}
\newcommand\field{\Bbbk}
\newcommand\fie{\Bbbk}
\def\Sspan#1{\langle #1\rangle}
\def\Span#1{\mathrm{span}(#1)}
\def\alg{\mathbb{A}}
\def\Alg{\mathbb{A}}
\def\balg{\mathbb{B}}
\def\valg{\mathbb{V}}
\def\ualg{\mathbb{U}}
\def\walg{\mathbb{W}}
\def\halg{\mathbb{H}}
\def\ealg{\mathbb{E}}
\def\lalg{\mathbb{L}}
\def\salg{\mathbb{S}}
\def\yalg{Y}
\def\algy{\mathbb{Y}}
\def\congp{\stackrel{\centerdot}{\cong}}
\def\xp{x_0}
\def\xm{x_1}
\def\yp{y_0}
\def\ym{y_1}
\def\ve{\tilde{e}}
\def\vt{{v}}
\def\xt{{x}}
\def\wt{{w}}
\newcommand{\fusion}{\mathscr{F}}
\newcommand{\basis}{\mathscr{B}}
\newcommand{\mult}{\cdot}
\newcommand{\plusalt}{\ast}
\newcommand{\Id}{\operatorname{Id}}
\newcommand{\Al}{\mathcal{A}}
\newcommand{\graf}[2]{\langle #1 | #2\rangle}
\newcommand{\nass}{\operatorname{\mathsf{NC}}}
\newcommand{\Nring}[2]{#1 (#2)}
\newcommand{\Part}{\operatorname{\mathbb{Part}}}
\newcommand{\dum}{\,\cdot\,\,}
\newcommand{\eval}{\operatorname{ev}}
\newcommand{\evan}{\mathcal{E}}
\DeclareMathOperator{\Tri}{\mathcal{S}}
\DeclareMathOperator{\Res}{\mathcal{R}}
\DeclareMathOperator{\Fix}{Fix}
\DeclareMathOperator{\interior}{int}
\DeclareMathOperator{\type}{type}
\DeclareMathOperator{\Div}{div} \DeclareMathOperator{\rank}{rank}
\DeclareMathOperator{\Dis}{Dis} \DeclareMathOperator{\re}{Re}
\DeclareMathOperator{\im}{im} \DeclareMathOperator{\Sec}{Sec(\V)}
\DeclareMathOperator{\Han}{Han}
\newcommand{\chr}{\operatorname{char}}
\DeclareMathOperator{\Char}{char} \DeclareMathOperator{\sing}{sing}
\DeclareMathOperator{\Idm}{Idm}
\DeclareMathOperator{\Nil}{Nil_2}
\DeclareMathOperator{\Cliff}{Cliff}
\DeclareMathOperator{\trace}{tr}
\DeclareMathOperator{\tr}{tr}
\DeclareMathOperator{\card}{card}
\DeclareMathOperator{\Affin}{Affin}
\DeclareMathOperator{\sgn}{sgn}
\DeclareMathOperator{\Aut}{Aut}
\DeclareMathOperator{\End}{End}
\DeclareMathOperator{\ord}{ord}
\DeclareMathOperator{\image}{im}
\DeclareMathOperator{\spec}{spec}
\DeclareMathOperator{\subs}{subs}
\DeclareMathOperator{\supp}{supp}
\DeclareMathOperator{\Spec}{Spec}
\newcommand{\NCA}{\operatorname{NCA}}
\DeclareMathOperator{\St}{St}
\DeclareMathOperator{\Ax}{Ax}
\DeclareMathOperator{\Sym}{Sym}
\DeclareMathOperator{\weight}{deg}
\DeclareMathOperator{\height}{ht}
\def\ualg{\mathbb{U}}
\def\malg{\mathbb{M}}
\newcommand{\vm}[1]{\malg(#1)}
\newcommand{\refl}{h}
\begin{document}
\title[On the Peirce spectral rigidity of decorated incidence algeb]{On the Peirce spectral rigidity of decorated incidence algebras}
\author{Daniel J. F. Fox; Vladimir G. Tkachev}
\date{}
\maketitle
\noindent\emph{Source and licence.} On the Peirce spectral rigidity of decorated incidence algebras. Version arXiv:2609.25353v1 (CC BY 4.0). This material is distributed under the Creative Commons Attribution 4.0 International licence, http://creativecommons.org/licenses/by/4.0/, as recorded in the arXiv OAI licence metadata for this exact version and shown on the abstract page https://arxiv.org/abs/2609.25353v1. Full rights notice: licenses/arxiv-2609.25353-CC-BY-4.0.txt.\par\smallskip
\noindent\textit{Editorial scope.} Counted unit: the original \texttt{theorem} environment \texttt{thm:localpaschspectrum} at source lines 960--1002 together with its complete proof at lines 1005--1047; the delivered document keeps source lines 432--1048 so that every referenced label and every citation resolves inside it. Native editable source is the paper's own concatenated TeX; the AMS wrapper is editorial formatting only. No publisher class, upstream script or unrelated artwork is redistributed.
\section{Spectral analysis of block idempotents}

\subsection{Preliminaries}\label{sec:not}

For a finite set $X$, $\binom{X}{k}$ denotes the set of $k$-element subsets of $X$. Fix an integer $m\ge2$. Throughout the paper $\Omega:=\{1,\dots,m\}$.
The Grassmannian geometry $G_2(m)$ is the partial Steiner triple system whose point set is
$\mathcal P=\binom{\Omega}{2}$,
and whose set of blocks $\mathcal B$ comprises the triples of pairwise collinear points
$b_{ijk}:=\{ij,ik,jk\}$, which are parametrized by $3$-element subset $\{i,j,k\}\in\binom{\Omega}{3}$. Their cardinalities are $|\mathcal P|=\binom{m}{2}$ and $|\mathcal B|=\binom{m}{3}$.
Similarly, quadrangles (Pasch configurations) of $G_2(m)$ are parametrized by $\binom{\Omega}{4}$. We shall routinely identify the point $\{i,j\}\in\binom{\Omega}{2}$ with the symbol $ij$.

The \emph{support} of the block $b\in\mathcal B$,
\begin{equation}\label{eq:suppdef}
\supp(b):=\bigcup_{p\in b}p\ \subset\ \Omega,
\end{equation}
is the triple of the elements of $\Omega$ contained in the points of $b$, so that $\supp(b_{ijk}) = \{i,j,k\}$. The map $b\mapsto\supp(b)$ is a bijection between $\mathcal B$ and $\binom{\Omega}{3}$, with inverse $\{i,j,k\}\mapsto b_{ijk} = \{ij,ik,jk\}$.







Next we describe the decorated incidence algebras associated with the combinatorial geometry $G_2(m)$. Given a \textit{decoration} $\delta:\mathcal B\to\{\pm1\}$,
the \textit{incidence algebra} $\alg^\delta(G_2(m))$ is the vector space
$\valg=\bigoplus_{p\in\mathcal P}\fie e_p$ with basis
$\{e_p:p\in\mathcal P\}$ equipped with the commutative multiplication given by extending bilinearly the products
\begin{equation}\label{multiplication}
e_{ij}\circ_\delta e_{ik}
=
\delta(\{ij,ik,jk\})\,e_{jk},
\qquad
\{i,j,k\}\in\binom{\Omega}{3},
\end{equation}
together with
\begin{align}\label{multiplication2}
&e_p\circ_\delta e_q=0,& &p, q \in \mathcal P \text{ not collinear.}
\end{align}
In particular, $e_p\circ_\delta e_p = 0$ so the generators $e_p$ are $2$-nilpotents.

%Whenever dependence on the decoration $\delta:\mathcal B\to\{\pm1\}$ is relevant, it will be indicated with a subscript $\delta$. Otherwise, to simplify notation, this dependence will be suppressed notationally.




The following definition from \cite{Fox-Tka2026b} plays a central role in the spectral theory of decorated partial Steiner triple systems.

\begin{definition}[\cite{Fox-Tka2026b}]\label{def:paschsign}
Let $\Pi=(\mathcal P,\mathcal B)$ be a decorated partial Steiner triple system with decoration
$\delta:\mathcal B\to\{\pm1\}.$
If $P=\{b_1,b_2,b_3,b_4\}\subseteq\mathcal B$
is a Pasch configuration, then its \emph{$\delta$-sign} is defined by
\begin{equation}\label{kappaP}
\kappa_\delta(P)
=
\prod_{b\in P}\delta(b)
=
\delta(b_1)\delta(b_2)\delta(b_3)\delta(b_4).
\end{equation}
\end{definition}

In the Grassmannian geometry $G_2(m)$, a Pasch configuration is uniquely determined by a $4$-subset
$\rho = \{i,j,k,r\}\in\binom{\Omega}{4}$. Precisely, $\rho$ determines the Pasch configuration
$$
P_\rho
=
\{
b_{ijk},
b_{ijr},
b_{ikr},
b_{jkr}
\}.
$$
Accordingly, it is convenient to write
\begin{equation}\label{kappaijkr}
\kappa_{ijkr}
:=
\kappa_\delta(P_\rho)
=
\delta(b_{ijk})
\delta(b_{ijr})
\delta(b_{ikr})
\delta(b_{jkr}).
\end{equation}




\subsection{A block decomposition of $G_2(m)$}\label{sec:proof}
Fix a block $b\in\mathcal B$, write $\{i,j,k\}:=\supp(b)$ so that $b=\{ij,ik,jk\}$, and define the cardinality $m-3$ subset $R:=\Omega\setminus\{i,j,k\}$. The underlying point set of $G_2(m)$ decomposes naturally as disjoint union
\begin{equation}\label{Pdecomp}
\mathcal{P}
=
\mathcal{P}_0
\sqcup
\mathcal{P}_1
\sqcup
\mathcal{P}_2,
\end{equation}
of the subsets
\begin{align*}
&\mathcal{P}_0=
b
=
\{ij,ik,jk\},&
&\mathcal{P}_1=
\bigsqcup_{r\in R}
\{ir,jr,kr\},
&&
\mathcal{P}_2=
\{rs:r,s\in R,\ r\ne s\},
\end{align*}
having cardinalities
\begin{align*}
&|\mathcal{P}_0|=3,&&
|\mathcal{P}_1|=3(m-3),&&
|\mathcal{P}_2|=\binom{m-3}{2}.
\end{align*}
For each $r\in R$, define
$$
W_r:=\{ir,jr,kr\},
\qquad
r\in R.
$$
As $\mathcal{P}_2 =
\binom{R}{2}$ consists of all points completely outside the fixed block, it is identified with the point set of the Grassmannian
$G_2(m-3)$ based on the $(m-3)$-set $R$. Thus the decomposition \eqref{Pdecomp} can be rewritten slightly imprecisely as
$$
\mathcal{P}
=
b
\sqcup
\Bigl(
\bigsqcup_{r\in R}W_r
\Bigr)
\sqcup
G_2(m-3).
$$
Passing from geometry to algebra there results a corresponding decomposition
\begin{equation}\label{Vdecomp}
\valg
=
\valg_{\mathcal{P}_0}
\oplus
\Bigl(
\bigoplus_{r\in R}\valg_{W_r}
\Bigr)
\oplus
\valg_{\mathcal{P}_2},
\end{equation}
where
$\valg_S
:=
\fie\{ e_p:\ p\in S\}$
denotes the linear span of $S$.


\subsection{A cohomological complex and a canonical $\mathbb Z_2$-grading}\label{sec:d-complex}

Both gauge transformations and Pasch-sign patterns admit a natural cohomological
interpretation. For Grassmannian geometries the relevant simplicial
complex is the full simplex on the vertex set $\Omega$: its
$1$-simplices are the points of $G_2(m)$, its $2$-simplices are the
blocks, and its $3$-simplices are the Pasch configurations.

Identifying $\{\pm1\}$ with $\mathbb F_2$, vertex signs
$\eta:\Omega\to\{\pm1\}$, point sign changes $\varepsilon$,
decorations $\delta$, and Pasch-sign patterns $\kappa_\delta$
can be regarded as $0$-, $1$-, $2$-, and $3$-cochains, respectively. The corresponding
coboundary operators
\begin{equation}\label{ddef2}
C^0\stackrel{d_0}{\longrightarrow}
C^1\stackrel{d_1}{\longrightarrow}
C^2\stackrel{d_2}{\longrightarrow}
C^3 .
\end{equation}
are given by
\begin{align}\label{d0d1d2}
\begin{aligned}
(d_0\eta)(ij)
&=
\eta(i)\eta(j),\\
(d_1\varepsilon)(ijk)
&=
\varepsilon(ij)\varepsilon(jk)\varepsilon(ki),\\
(d_2\delta)(ijkl)
&=
\delta(ijk)\delta(ijl)\delta(ikl)\delta(jkl),
\end{aligned}
\end{align}
and coincide with the assignments
\begin{equation}\label{ddef1}
\eta\stackrel{d_0}{\longmapsto}\varepsilon,
\qquad
\varepsilon\stackrel{d_1}{\longmapsto}\delta^\varepsilon/\delta,
\qquad
\delta\stackrel{d_2}{\longmapsto}\kappa_\delta.
\end{equation}
Indeed, the three points of the block $\{i,j,k\}\in\binom{\Omega}{3}$ are
$ij$, $jk$ and $ki$, so that
$$
\delta^\varepsilon(\{i,j,k\})
=
\varepsilon(ij)\varepsilon(jk)\varepsilon(ki)\,
\delta(\{i,j,k\})
=
(d_1\varepsilon)(\{i,j,k\})\,\delta(\{i,j,k\}),
$$
that is, the gauge action is translation of the decoration space by the
subgroup $\operatorname{im}d_1$, while the Pasch-sign pattern is obtained
from the decoration by the next coboundary operator,
$\kappa_\delta=d_2\delta$. In particular, the gauge invariance of the
Pasch signs recorded above is nothing but the identity
\begin{equation}\label{d2d1eq}
d_2\circ d_1=0.
\end{equation}
Since the full simplex is contractible, this complex is exact in positive
degrees. Consequently,
\begin{equation}\label{rankd2}
\ker d_2=\operatorname{im}d_1,
\end{equation}
so that two decorations of $G_2(m)$ have the same Pasch-sign pattern if
and only if they are gauge equivalent; in particular, the realizable
Pasch-sign patterns are in bijection with the gauge-equivalence classes.
Furthermore,
$$
\ker d_1=\operatorname{im}d_0,
$$
so the gauge kernel consists exactly of the sign changes
$$
\varepsilon(ij)=\eta(i)\eta(j),
\qquad
i,j\in\Omega,
$$
and is therefore isomorphic to $\mathbb F_2^{m-1}$. It follows that
\begin{equation}\label{rankd1}
\operatorname{rank}d_1
=
\binom{m}{2}-(m-1)
=
\binom{m-1}{2}.
\end{equation}



\begin{lemma}\label{lem:gradingauto}
For a subset $\tau\subseteq\Omega$ define
$\varepsilon_\tau:\mathcal P\to\{\pm1\}$ by
\begin{equation}\label{epstau}
\varepsilon_\tau(p)=(-1)^{|p\cap\tau|},
\qquad
p\in\mathcal P,
\end{equation}
that is, $\varepsilon_\tau=d_0(\mathbf 1_\tau)$, where
$\mathbf 1_\tau\in C^0$ is the indicator function of $\tau$. Then, for
\emph{every} decoration $\delta$, the diagonal map
\begin{equation}\label{thetatau}
\theta_\tau:\ \alg^\delta(G_2(m))\to\alg^\delta(G_2(m)),
\qquad
\theta_\tau(e_p)=\varepsilon_\tau(p)\,e_p ,
\end{equation}
is an involutive automorphism, and the induced decomposition
\begin{equation}\label{taugrading}
\begin{aligned}
\alg^\delta&=\alg_{+}(\tau)\oplus\alg_{-}(\tau),\\
\alg_{\pm}(\tau)&=\valg_{\mathcal P^{\pm}(\tau)},\\
\mathcal P^{\pm}(\tau)&=\{p\in\mathcal P:\ (-1)^{|p\cap\tau|}=\pm1\},
\end{aligned}
\end{equation}
is a $\mathbb Z_2$-grading of $\alg^\delta$.
\end{lemma}

\begin{proof}
Clearly $\theta_\tau^2=\mathrm{id}$. That $\theta_{\tau}$ is a homomorphism need only be tested
on basis vectors. If $p,q$ are not collinear then
$e_p\circ_\delta e_q=0$, so
$$\theta_\tau(e_p)\circ_\delta\theta_\tau(e_q)
= \epsilon_{\tau}(p)\epsilon_{\tau}(q)e_{p}\circ e_{q} = 0 = \theta_{\tau}(e_{p}\circ_{\delta}e_{q}).$$
If $p,q,r$ are the three points of a block $b'$,
then $e_p\circ_\delta e_q=\delta(b')e_r$, so that
$$
\theta_\tau(e_p)\circ_\delta\theta_\tau(e_q)
=\varepsilon_\tau(p)\varepsilon_\tau(q)\,\delta(b')e_r,
\qquad
\theta_\tau(e_p\circ_\delta e_q)=\varepsilon_\tau(r)\,\delta(b')e_r .
$$
These agree for every block $b'$ if and only if
$\varepsilon_\tau(p)\varepsilon_\tau(q)\varepsilon_\tau(r)=1$, that is,
if and only if $d_1\varepsilon_\tau\equiv1$, which holds because
$d_1\varepsilon_\tau=d_1d_0(\mathbf 1_\tau)=0$. This proves that
$\theta_\tau$ is an automorphism, and \eqref{taugrading} is its
eigenspace decomposition by \eqref{epstau}.
\end{proof}

Specializing $\tau=\supp{b}$ yields, for every block $b$, a canonical grading automorphism of the incidence
algebra independently of the decoration $\delta$.

\begin{corollary}[A canonical $\mathbb Z_2$-grading]
\label{cor:gradingauto}
Let $b\in\mathcal B$ and put
$\varepsilon_b:=\varepsilon_{\supp(b)}$, $\theta_b:=\theta_{\supp(b)}$,
so that
\begin{equation}\label{epsilonmap}
\varepsilon_b(p)=(-1)^{|p\,\cap\,\supp(b)|},
\qquad
\theta_b(e_p)=\varepsilon_b(p)\,e_p .
\end{equation}
Then, for every decoration $\delta$, the map $\theta_b$ is an involutive
automorphism of $\alg^\delta(G_2(m))$, and its eigenspace decomposition
is
\begin{align}\label{apriorigrading}
\begin{aligned}
&\alg^\delta=\alg_{+}(b)\oplus\alg_{-}(b),\\
&\alg_{+}(b)=\valg_{\mathcal P_0}\oplus\valg_{\mathcal P_2},&&
&\alg_{-}(b)=\valg_{\mathcal P_1}=\bigoplus_{r\in R}\valg_{W_r} .
\end{aligned}
\end{align}
\end{corollary}

\begin{proof}
Apply Lemma~\ref{lem:gradingauto} with $\tau=\supp(b)$. By
\eqref{Pdecomp}, $|p\cap\supp(b)|$ equals $2$ for $p\in\mathcal P_0$,
$1$ for $p\in\mathcal P_1$ and $0$ for $p\in\mathcal P_2$; hence
$\mathcal P^{+}(\tau)=\mathcal P_0\sqcup\mathcal P_2$ and
$\mathcal P^{-}(\tau)=\mathcal P_1$, which is \eqref{apriorigrading}.
\end{proof}


\subsection{The Peirce decomposition}
The decomposition \eqref{Vdecomp} is  a consequence of the fact that the multiplication in $\alg^\delta(G_2(m))$ is compatible with the natural inclusions of smaller Grassmannian geometries. The following lemma makes this observation precise.

\begin{lemma}\label{lem:closure}
For $T\subseteq\Omega=\{1,\dots,m\}$ the subspace
$\valg(T):=\fie\{e_{ij}: i,j\in T\}$ is a subalgebra of $\alg^\delta(G_2(m))$ having dimension $
\dim \valg(T)=\binom{|T|}{2}$, and such that if $\delta_T$ denotes the restriction of $\delta$ to the blocks contained in $T$, then $\valg(T)$ is isomorphic with $\alg^{\delta_T}(G_2(|T|))$,
$\valg(T)\simeq\alg^{\delta_T}(G_2(|T|))$.
\end{lemma}

\begin{proof}
If $p,q\in\mathcal P(T)$ are not collinear, then $e_p\circ_\delta e_q=0$. If $p=ij$ and $q=ik$, then \[ e_{ij}\circ_\delta e_{ik} = \delta(b_{ijk})e_{jk}, \] and since $j,k\in T$, the product again belongs to $\valg(T)$. Hence $\valg(T)$ is a subalgebra. The remaining assertions are immediate.
\end{proof}



\begin{corollary}\label{cor:where} Let $S\subseteq\mathcal{P}$ and let $$ T:=\bigcup_{ij\in S}\{i,j\}. $$ Then the subalgebra generated by $\{e_p:p\in S\}$ is contained in $\valg(T).$ \end{corollary}

The decomposition \eqref{Vdecomp} is adapted to the geometry of the fixed block and serves as the starting point for the spectral analysis. Proposition \ref{prop:blockdecomp} shows that it is preserved by multiplication with the block idempotent.

\begin{proposition}\label{prop:blockdecomp}
For any decoration $\delta:\mathcal{B}\to\{\pm1\}$ and any block $b\in\mathcal B$, the decomposition \eqref{Vdecomp} is invariant under $L_\delta(c_b^\delta)$. Moreover, with respect to \eqref{Vdecomp}, the operator $L_\delta(c_b^\delta)$ is block diagonal,
\begin{equation}\label{blockdiag}
L_\delta(c_b^\delta)
=
L_0
\oplus
\Bigl(
\bigoplus_{r\in R}L_r
\Bigr)
\oplus
0_{\mathcal{P}_2}.
\end{equation}
Here $L_0$ is the restriction of $L_\delta(c_b^\delta)$ to $\valg_{\mathcal{P}_0}$, $L_r$ is its restriction to $\valg_{W_r}$, and $0_{\mathcal{P}_2}$ denotes the zero operator on $\valg_{\mathcal{P}_2}$.
\end{proposition}

\begin{proof}
Write $c=c_b^\delta.$ Since $c$ is supported on
$$
\mathcal{P}_0=b=\{ij,ik,jk\},\quad \{i,j,k\}=\supp(b),
$$
it suffices to determine the action of $c$ on the three summands of \eqref{Vdecomp}.
Since $\mathcal{P}_0$ is a block, by Corollary~\ref{cor:where}, $\valg_{\mathcal{P}_0}$ is a subalgebra containing $c$, so
\begin{equation}\label{B0inv}
L_\delta(c)(\valg_{\mathcal{P}_0})
\subseteq
\valg_{\mathcal{P}_0}.
\end{equation}
For $W_r=\{ir,jr,kr\}$, $r\in R$, there hold
\begin{align}\label{wr3}
&e_{ir}\circ_\delta e_{ij}=\delta(b_{ijr}) e_{jr},&
&e_{ir}\circ_\delta e_{ik}=\delta(b_{ikr}) e_{kr},&
&e_{ir}\circ_\delta e_{jk}=0,
\end{align}
and similar identities hold after cyclic permutation of $i,j,k$. Therefore
\begin{equation}\label{wrinv}
L_\delta(c)(\valg_{W_r})
\subseteq
\valg_{W_r}
\end{equation}
for every $r\in R$. Finally, if $rs\in\mathcal{P}_2$, then $rs$ does not belong to any block containing a point of $\mathcal{P}_0$, so
\begin{align}\label{B2a}
&e_{rs}\circ_\delta e_{ij}=0,
&&
e_{rs}\circ_\delta e_{ik}=0,
&&
e_{rs}\circ_\delta e_{jk}=0,
\end{align}
and consequently
\begin{equation}\label{B2zero}
L_\delta(c)(\valg_{\mathcal{P}_2})=0.
\end{equation}
Combining \eqref{B0inv}, \eqref{wrinv}, and \eqref{B2zero} yields the block decomposition \eqref{blockdiag}.
\end{proof}

The block decomposition reduces the spectral problem to the computation of the individual diagonal blocks. We now determine these blocks explicitly.

\begin{proposition}\label{prop:blockmatrices}
Let $b\in\mathcal B$, write $\{i,j,k\}:=\supp(b)$, and let $c=c_b^\delta$. The matrix of $L_0$ with respect to the basis $\{e_{ij},e_{ik},e_{jk}\}$ is
\begin{equation}\label{L0matrix}
L_0
=
\tfrac12
\begin{pmatrix}
0&1&1\\
1&0&1\\
1&1&0
\end{pmatrix},
\end{equation}
and therefore
\begin{equation}\label{specL0}
\Spec(L_0)
=
\Bigl\{
1,
\Bigl(-\tfrac12\Bigr)^2
\Bigr\}.
\end{equation}
Furthermore, for every $r\in R$, the characteristic polynomial of $L_r$ is
\begin{equation}\label{charpolyr}
\chi_r(\lambda)
=
\lambda^3-\tfrac34\lambda-\tfrac14\kappa_{ijkr}.
\end{equation}
Consequently,
\begin{equation}\label{specplus}
\Spec(L_r)
=
\left\{
  \begin{array}{ll}
    \bigl\{1,\bigl(-\tfrac12\bigr)^2\bigr\}, & \hbox{if $\kappa_{ijkr}=1$;} \\
    \bigl\{-1,\bigl(\tfrac12\bigr)^2\bigr\}, & \hbox{if $\kappa_{ijkr}=-1$.}
  \end{array}
\right.
\end{equation}
In particular,
\begin{equation}\label{fiveeigen}
\spec\bigl(L_\delta(c_b^\delta)\bigr)
\subseteq
\Bigl\{
0,\pm1,\pm\tfrac12
\Bigr\}.
\end{equation}
\end{proposition}


\begin{proof}
By Proposition \ref{prop:blockdecomp}, the operator $L_\delta(c)$ preserves each summand in \eqref{Vdecomp}, acts on $\valg_{\mathcal{P}_2}$ as zero, and decomposes as in \eqref{blockdiag}. It remains only to compute the non-zero blocks. We compute separately the blocks $L_0$ and $L_r$.
First consider the summand $\valg_{\mathcal{P}_0}$, where $\mathcal{P}_0=\{ij,ik,jk\}$.
By definition of the block idempotent $c_b^\delta$, using $e_{ij}\circ_\delta e_{ij}=0$ and $\delta(b)^2=1$, a direct computation gives
$$
L_\delta(c)e_{ij}
=
\tfrac12\delta(b)\bigl(e_{ij}\circ_\delta e_{ij}+e_{ik}\circ_\delta e_{ij}+e_{jk}\circ_\delta e_{ij}\bigr)
=
\tfrac12\bigl(e_{jk}+e_{ik}\bigr),
$$
and similarly for $e_{ik}$ and $e_{jk}$. Hence the restriction of $L_\delta(c)$ to $\valg_{\mathcal{P}_0}$ is represented by the matrix $L_{0}$ in \eqref{L0matrix} which has characteristic polynomial
$$
\chi_0(\lambda)
=
\lambda^3-\tfrac34\lambda-\tfrac14 =
(\lambda-1)\Bigl(\lambda+\tfrac12\Bigr)^2.
$$
This proves \eqref{specL0}.
Now fix $r\in R$ and consider $W_r=\{ir,jr,kr\}$. The corresponding Pasch configuration $
P:=\{b_{ijk},
b_{ijr},
b_{ikr},
b_{jkr}\}
$ has $\delta$-sign $\kappa_{ijkr}$ as in \eqref{kappaijkr}.
In the ordered basis $(e_{ir},e_{jr},e_{kr}),$
the operator $L_r$ is represented by a matrix of the form
$$
L_r
=
\frac12
\begin{pmatrix}
0&s_1&s_2\\
s_1&0&s_3\\
s_2&s_3&0
\end{pmatrix},
$$
where
\begin{align}
s_1&=\delta(b_{ijk})\delta(b_{ijr}),
\quad
s_2=\delta(b_{ijk})\delta(b_{ikr}),
\quad
s_3=\delta(b_{ijk})\delta(b_{jkr}).
\label{eps3}
\end{align}
Since $\delta(b_{ijk})^2=1$, it follows that
\begin{equation}\label{epsproduct}
s_1s_2s_3
=
\delta(b_{ijk})
\delta(b_{ijr})
\delta(b_{ikr})
\delta(b_{jkr})
=
\kappa_{ijkr}.
\end{equation}


Using \eqref{epsproduct}, a  direct determinant calculation yields
$$
\begin{aligned}
\chi_r(\lambda)&=\det(\lambda I-L_r) =\lambda^3-\tfrac34\lambda-\tfrac14 s_1s_2s_3
\\&= \lambda^3-\tfrac34\lambda-\tfrac14\kappa_{ijkr}  = (\lambda-\kappa_{ijkr})\bigl(\lambda+\tfrac12 \kappa_{ijkr}\bigr)^2,
\end{aligned}
$$
which shows \eqref{charpolyr} and \eqref{specplus}.
Finally, by \eqref{blockdiag}, the remaining block is $0_{\mathcal{P}_2}$. Combining \eqref{specL0}, \eqref{specplus} and \eqref{B2zero} yields \eqref{fiveeigen}.
\end{proof}

\subsection{Proof of the main result}\label{sec:localpasch}

Combining Propositions~\ref{prop:blockdecomp} and \ref{prop:blockmatrices}, we obtain the main result of the paper, which expresses the complete multiset spectrum of a block idempotent in terms of the number of negative local Pasch configurations containing the corresponding block.


\begin{theorem}\label{thm:localpaschspectrum}
Let $b\in\mathcal B$, write $\{i,j,k\}:=\supp(b)$, let $R:=\Omega\setminus\{i,j,k\}$, and let
\begin{equation}\label{kjkr}
\kappa_{ijkr}
=
\delta(b)
\delta(b_{ijr})
\delta(b_{ikr})
\delta(b_{jkr}),
\qquad
r\in R.
\end{equation}
Define
\begin{equation}\label{nminusone1}
n^{\delta}_{-1}(b)
=
|\{r\in R:\kappa_{ijkr}=-1\}|,
\end{equation}
which we abbreviate to $n^{\delta}_{-1}$ whenever the block $b$ is
fixed. Then
\begin{equation}\label{nplusrelation}
|\{r\in R:\kappa_{ijkr}=1\}| = m-3-n^{\delta}_{-1}
\end{equation}
and
\begin{equation}\label{localspectrumformula}
\begin{aligned}
\Spec&(L_\delta(c_b^\delta))=\sigma_{n^{\delta}_{-1}}
%&=\Bigl\{ 0^{\binom{m-3}{2}}, 1^{m-2-n^{\delta}_{-1}}, (-1)^{n^{\delta}_{-1}}, (-\tfrac12)^{2(m-2-n^{\delta}_{-1})}, (\tfrac12)^{2n^{\delta}_{-1}} \Bigr\}.
\end{aligned}
\end{equation}
where
\begin{equation}\label{sigmanot}
\sigma_k
=
\Bigl\{
0^{\binom{m-3}{2}},
1^{m-2-k},
(-1)^k,
(-\tfrac12)^{2(m-2-k)},
(\tfrac12)^{2k}
\Bigr\}
\end{equation}
\end{theorem}

\begin{proof}
Put $c:=c_b^\delta.$
Since $|R|=m-3,$
it follows from the definition of $n^{\delta}_{-1}$ that
\begin{equation}\label{partitionR}
R=
\{r\in R:\kappa_{ijkr}=1\}
\sqcup
\{r\in R:\kappa_{ijkr}=-1\}.
\end{equation}
Taking cardinalities in \eqref{partitionR} yields
\begin{equation}\label{nplusrelationproof}
|\{r\in R:\kappa_{ijkr}=1\}|
=
m-3-n^{\delta}_{-1},
\end{equation}
which proves \eqref{nplusrelation}.
By Proposition \ref{prop:blockdecomp}, $L_\delta(c)$ preserves the decomposition \eqref{Vdecomp} with respect to which it has the block form \eqref{blockdiag}.
Therefore
$\Spec(L_\delta(c))$
is obtained by combining the spectra of the diagonal blocks. By Proposition \ref{prop:blockmatrices}, the spectra of the blocks $L_0$ and $L_r$ are given in \eqref{specL0} and \eqref{specplus}.
By \eqref{nplusrelationproof}, there are $m-3-n^{\delta}_{-1}$ elements $r\in R$ satisfying $\kappa_{ijkr}=1$, and $n^{\delta}_{-1}$ elements $r\in R$ satisfying $\kappa_{ijkr}=-1$.
Hence the blocks $L_r$ contribute
\begin{equation}\label{Wrcontribution}
\Bigl\{
1^{m-3-n^{\delta}_{-1}},
(-1)^{n^{\delta}_{-1}},
(-\tfrac12)^{2(m-3-n^{\delta}_{-1})},
(\tfrac12)^{2n^{\delta}_{-1}}
\Bigr\}
\end{equation}
to the multiset spectrum.
The only remaining contribution comes from the residual Grassmannian component. Proposition \ref{prop:blockdecomp} implies that $L_\delta(c)$ vanishes on $\valg_{\mathcal{P}_2}$,
and since
$
\dim\valg_{\mathcal{P}_2}
=
|\mathcal{P}_2|
=
\binom{m-3}{2},
$
this block contributes $0^{\binom{m-3}{2}}.$ Combining this with \eqref{specL0} and \eqref{Wrcontribution} yields \eqref{localspectrumformula}.
\end{proof}

\begin{thebibliography}{99}
\bibitem[Fox-Tka2026b]{Fox-Tka2026b} D.~J.~F. Fox and V.~G. Tkachev, \emph{Spectral rigidity and {P}eirce decompositions of decorated partial {S}teiner triple systems}, 2026.
\end{thebibliography}
\end{document}
